% =============================================================
% Ucapan Terima Kasih (Bab 3.9 of the Pedoman).
% Keep this concise; thank persons/institutions that helped
% materially or non-materially.
% =============================================================
\thispagestyle{empty}
\begin{center}
{\bfseries UCAPAN TERIMA KASIH\par}
\vspace{10mm}
\end{center}

\noindent Puji syukur penulis panjatkan kepada Tuhan Yang Maha Esa atas rahmat dan karunia-Nya sehingga tesis yang berjudul \textit{``Judul Tesis untuk Program Magister (S2) Data Science dan Kecerdasan Buatan''} ini dapat diselesaikan sebagai salah satu syarat memperoleh gelar Magister pada Program Studi S2 Data Science dan Kecerdasan Buatan, Fakultas Ilmu Komputer dan Teknologi Informasi, Universitas Sumatera Utara.

Penulis mengucapkan terima kasih kepada Komisi Pembimbing atas bimbingan, arahan, dan waktu yang telah diberikan selama proses penelitian dan penulisan tesis ini. Ucapan terima kasih juga disampaikan kepada seluruh dosen dan staf Program Studi S2 Data Science dan Kecerdasan Buatan, keluarga, serta rekan-rekan yang telah memberikan dukungan moril maupun materiil.

Penulis menyadari bahwa tesis ini masih memiliki kekurangan. Oleh karena itu, kritik dan saran yang membangun sangat diharapkan demi kesempurnaan penelitian selanjutnya.

\vspace{10mm}
\noindent Medan, \the\year

\vspace{6mm}
\noindent Penulis
\clearpage
